\section{Discussão}\label{sec:Discussion}

\subsection{Baixa previsibilidade pré-jogo}\label{subsec:DiscPrevisibilidade}

O resultado central é que o histórico das equipes disponível antes da partida explica pouco do total de gols e de escanteios. O melhor modelo reduziu a desviância da média da liga em cerca de 1,7\% nos gols e de 0,9\% nos escanteios, e a variação da desviância entre temporadas, mesmo para a média da liga, é maior que a diferença entre os modelos. A redução dos gols resiste a referências mais fortes que a média da liga, mas a dos escanteios não: contra uma média que acompanha o nível de cada temporada, nenhum modelo tem ganho significativo (Seção~\ref{subsec:ResReferencias}). A dimensão do sinal aparece também na correlação entre a previsão do Random Forest e o total observado, calculada dentro de cada temporada de teste: a média das seis correlações é de 0,14 nos gols e de 0,09 nos escanteios. A perda logarítmica das probabilidades de passar das linhas de 2,5 gols e de 9,5 escanteios ficou próxima da de uma previsão de 50\% para todas as partidas. Isso não indica falha dos modelos, mas o limite da informação usada: o total de uma partida depende de eventos raros e de circunstâncias do próprio jogo, como um gol cedo, uma expulsão ou a necessidade de reagir ao placar, que nenhuma variável pré-jogo antecipa.

Três observações reforçam essa leitura. Primeiro, as árvores escolhidas na validação interna foram rasas, muitas vezes com uma única divisão ou com folhas grandes, o que indica que o sinal disponível é fraco e quase aditivo, sem interações que modelos mais flexíveis pudessem explorar. Segundo, o XGBoost, o modelo mais flexível, não superou o Random Forest, e nenhuma combinação das previsões superou o melhor modelo isolado com significância. Terceiro, as escolhas de configuração variaram entre dobras vizinhas, o que é esperado quando as alternativas diferem pouco entre si e a seleção é dominada pelo ruído. Esse quadro é coerente com a conclusão de \citeonline{berrar2019}, para quem o desempenho depende mais de como o conhecimento do domínio entra nas variáveis do que do algoritmo.

\subsection{Gols e escanteios}\label{subsec:DiscAlvos}

Os dois alvos se comportaram de forma diferente em três aspectos. O primeiro é a dispersão: os gols seguiram de perto a distribuição de Poisson, e a binomial negativa não trouxe ganho, enquanto os escanteios mostraram sobredispersão, e a binomial negativa melhorou as probabilidades em todos os modelos. Esse resultado é compatível com a observação de \citeonline{yip2024} de que os escanteios tendem a ocorrer em sequências dentro de uma mesma partida, o que aumenta a variância em relação à Poisson. Para a previsão da média, a distribuição pouco importa; para a previsão de uma linha, ignorar a sobredispersão dos escanteios piora as probabilidades.

O segundo aspecto é o tipo de informação útil. Nos gols, a contribuição veio dos índices de força, os \textit{pi-ratings} e os índices GAP, que resumem o confronto entre as equipes; nos escanteios, veio sobretudo da produção de escanteios de cada equipe, pela média da temporada, pela forma recente e pelos índices GAP de escanteios. Os escanteios parecem depender mais do estilo de jogo de cada equipe, que se mantém ao longo da temporada, do que da diferença de força entre elas.

O terceiro aspecto é a formação tática. A hipótese de que formações com mais amplitude pelos lados gerariam mais escanteios não se confirmou: a formação raramente foi escolhida na validação interna e, quando entrou, teve importância nula. Uma explicação possível é que o efeito da formação já esteja contido no histórico de escanteios da equipe, que tende a repetir o mesmo desenho tático; outra é que a formação inicial registrada descreva mal o posicionamento ao longo do jogo.

Embora os escanteios sejam mais frequentes que os gols, o ganho relativo dos modelos foi menor nos escanteios, e deixou de ser significativo contra a média corrente da liga. A maior frequência reduz a variância relativa da contagem, mas não torna o total mais previsível a partir do histórico: o que varia de uma partida para outra continua a depender do desenrolar do jogo.

\subsection{Comparação com os trabalhos relacionados}\label{subsec:DiscRelacionados}

A adequação da Poisson aos gols confirma, no Campeonato Brasileiro, o que \citeonline{maher1982} e \citeonline{dixoneColes1997} observaram no futebol inglês, e contrasta com a sobredispersão considerada por \citeonline{santana2019} e com as variações de Poisson comparadas por \citeonline{ribeiro2025}; esses dois estudos, porém, avaliaram uma única temporada da Série A e modelaram os gols de cada equipe, e não o total da partida. O bom desempenho do Random Forest contrasta com o de \citeonline{santana2019}, em que esse método teve desempenho ruim na previsão do resultado da Série A de 2017; a diferença de alvo, um total de contagem aqui e uma classe de resultado lá, e o uso de índices de força como variáveis podem explicar a diferença. O papel desses índices aproxima este trabalho de \citeonline{groll2019}, cujo Random Forest se beneficiou de estimativas de força das equipes obtidas por modelos de Poisson, e de \citeonline{hubacek2019} e \citeonline{yeung2024}, que também construíram variáveis a partir de índices de força, entre eles os \textit{pi-ratings}.

Os índices GAP foram o grupo mais importante nos gols, e tiveram papel menor nos escanteios, atrás da média da temporada e da forma recente, em linha com \citeonline{wheatcroft2020}, que os propôs justamente para totais de gols. Aqui, cada alvo usou apenas os índices da própria medida, e não se testou a combinação de chutes e escanteios para prever gols, que \citeonline{wheatcroft2020} encontrou ser mais informativa. A separação entre a estatística da partida-alvo e a sua previsão a partir do histórico, adotada por \citeonline{wheatcroft2021}, foi mantida em todas as variáveis, e a validação respeitou a ordem temporal das partidas, como recomendam \citeonline{tsokos2019}. Por fim, \citeonline{yip2024} obtiveram ganho nos escanteios ao incorporar as cotações de outros mercados de apostas, informação que este trabalho não usou; a ausência de ganho sobre a frequência na linha de 9,5 escanteios sugere que essa fonte externa pode ser necessária para prever escanteios melhor que a taxa histórica.

\subsection{Limitações}\label{subsec:DiscLimitacoes}

Além das limitações das fontes descritas na Seção~\ref{subsec:Limitacoes}, sete limitações decorrem da modelagem. Primeiro, a avaliação usa seis temporadas de teste, entre elas a de 2020, disputada sem público e fora do calendário habitual, e as diferenças entre modelos são da ordem da variação entre temporadas; os intervalos de confiança são largos em relação a essas diferenças. Segundo, várias comparações foram feitas sem correção para múltiplos testes, e o intervalo do Random Forest nos escanteios cruzou o zero por pouco, enquanto o da média das duas árvores não o cruzou; conclusões que dependem dessa margem são frágeis. A média das duas árvores, destaque dos escanteios, foi definida depois de vistos os resultados de teste de cada uma, e a significância dos escanteios mudou entre execuções (Tabela~\ref{tab:instabilidade}). O intervalo de confiança, por sua vez, mede só a incerteza da média sobre as partidas de teste com os modelos fixos (Seção~\ref{subsec:Modelos}). Terceiro, as grades do peso $k$ e do Random Forest foram ampliadas quando a escolha caiu na borda, mas o XGBoost explorou uma amostra de 240 das 15.552 combinações do seu espaço e ficou na borda da taxa de aprendizado em várias dobras. Para verificar se essa borda limitava o modelo, uma busca adicional acrescentou 240 combinações com taxa de aprendizado de 0,3, penalização L1 de 30 ou penalização L2 de 0,1. A desviância de validação das configurações escolhidas caiu no máximo 0,004, mas a de teste subiu de 1,0599 para 1,0617 nos gols e de 1,1751 para 1,1824 nos escanteios, alvo em que o XGBoost deixou de superar a média da liga. As novas escolhas caíram na borda de quase todos os hiperparâmetros, ora no extremo inferior, ora no superior, o que indica que a desviância de validação é quase plana nesse espaço e que ampliar a busca apenas ajusta o ruído das temporadas de validação. Os resultados relatados são os da busca original, declarada final antes da busca adicional, mas a decisão de mantê-la foi tomada depois de ver o teste da busca adicional, o que, na prática, é escolher entre duas execuções olhando o teste; com estes dados, ampliar a busca do XGBoost não se mostrou um caminho promissor. Quarto, a combinação ponderada foi avaliada com previsões de validação geradas com hiperparâmetros escolhidos nessas mesmas validações, o que a favorece um pouco; mesmo assim, ela não superou o melhor modelo. Quinto, a importância por permutação mede a dependência do modelo ajustado em cada variável, e não um efeito causal; no modelo linear, ela é distorcida pela correlação entre os índices de força, razão pela qual a discussão se apoiou nas árvores. Sexto, as dobras de 2019 e de 2020 escolhem a configuração com uma e com duas temporadas de validação, respectivamente, e a de 2019 valida em 2018 um modelo ajustado só em 2017, temporada de aquecimento, sem referência de temporada anterior; um terço do teste vem dessas duas dobras, em que a seleção é pouco confiável. Nos gols de 2019, o Poisson \textit{ridge} escolheu a penalização de $10^{-6}$, sem regularização prática, e ficou pior que a média da liga no teste (1,1546 contra 1,1314), e o $k$ escolhido nos gols variou entre 5, 1.000, 1.000, 5, 10 e 10 ao longo das dobras. Sétimo, não houve ablação por grupo de variáveis nem referência intermediária, como apenas o esperado do confronto pela média da temporada ou só os índices GAP; não se mediu, portanto, quanto da engenharia de variáveis foi necessária, e a importância por permutação, que sugere que nos gols quase tudo vem dos índices GAP e dos \textit{pi-ratings}, não substitui essa medida.

Ficaram fora informações que poderiam melhorar a previsão: desfalques e escalação nominal, cotações das casas de apostas, condições do gramado e do clima e a carga de partidas nos campeonatos estaduais e nas competições continentais. Os modelos também não foram comparados com as probabilidades implícitas nas cotações, de modo que não se pode afirmar se o ganho observado tem valor diante do mercado.

\subsection{Uso prático das linhas de 2,5 gols e 9,5 escanteios}\label{subsec:DiscLinhas}

Nas linhas usuais do mercado de apostas, os resultados indicam usos diferentes para os dois alvos. Para mais de 2,5 gols, o Random Forest e o XGBoost produziram probabilidades melhores que a frequência histórica, com significância, mas o ganho é pequeno: a perda logarítmica caiu cerca de 1\%. Como as cotações embutem a margem da casa e, possivelmente, informação que os modelos não têm, esse ganho não basta para afirmar que haveria lucro; a comparação com as cotações é o passo necessário para essa conclusão. Para mais de 9,5 escanteios, nenhuma previsão superou a frequência histórica, e a escolha da distribuição foi mais importante que a do modelo: a Poisson, ao subestimar a variância, piorou as probabilidades, e a binomial negativa apenas igualou a frequência. Na prática, a taxa histórica de partidas acima da linha é, com os dados usados, uma referência tão boa quanto os modelos para os escanteios, o que é coerente com a ausência de ganho significativo contra a média corrente da liga.
